\documentclass[UTF8, zihao=-4]{ctexart}
\usepackage{amsmath}
\usepackage{setspace}                    % 行距控制
\setstretch{1.0} % 严格单倍行距 (1.0倍)
\begin{document}

\section{问题三模型的建立与求解}
\subsection{数据分析与处理}

由分析，本小题第一问我们将用时间序列方法解答。首先进行数据准备与预处理.（如表1所示）
\begin{table}[ht]
    \centering
    \begin{tabular}{|c|c|}
    \hline
    \textbf{年份} & \textbf{全省教育经费投入（亿元）} \\
    \hline
    2015 & 597.52 \\
    2016 & 643.98 \\
    2017 & 658.67 \\
    2018 & 686.65 \\
    2019 & 677.74 \\
    2020 & 720.48 \\
    2021 & 709.2 \\
    2022 & 748.41 \\
    2023 & 835.33 \\
    \hline
    \end{tabular}
    \caption{吉林省教育经费投入情况（2015–2023年）}
    \label{tab:education_budget}
    \end{table}

我们以年份作为时间索引，以便对数据进行排序、索引和处理。由于数据没有缺失值，故不需要进行插值。
再将教育经费作为时序预测的数据列，以求未来三年的教育经费预测值。

考虑到近九年教育经费呈现明显的上升趋势（如下图），季节性变化较小，故我们选择双指数平滑模型对数据进行快速建模和预测。（这里插入上面表格的折线图，表现趋势。）

\subsection{建模求解}
我们使用双指数平滑模型将时间序列分解为两个成分：水平和趋势，并通过指数加权的方法对这两个成分进行平滑。

首先，我们给出水平更新公式、趋势更新公式及预测公式。

水平（Level）更新公式:
\[
\ell_t = \alpha y_t + (1 - \alpha)(\ell_{t-1} + b_{t-1})
\]

趋势（Trend）更新公式:
\[
b_t = \beta (\ell_t - \ell_{t-1}) + (1 - \beta) b_{t-1}
\]

预测公式:
\[
\hat{y}_{t+m} = \ell_t + m b_t
\]

其中，
\textbf{符号说明}：
\begin{itemize}
  \item $\ell_t$ 是时间点 $t$ 的水平估计；
  \item $b_t$ 是时间点 $t$ 的趋势估计；
  \item $y_t$ 是时间点 $t$ 的实际观测值，即本题中从2015年开始，每一年的教育经费。
  \item $\alpha$ 是水平平滑系数，控制对当前观测值的权重；
  \item $\beta$ 是趋势平滑系数，控制对趋势变化的敏感度；
  \item $h$ 是预测步长，取3,4,5，分别对应未来三年。
\end{itemize}

我们通过ExponentialSmoothing类，启用阻尼趋势，避免预测值无限增长。同时使用最小化均方误差法自动调整水平平滑系数和趋势平滑系数，最终寻找最优平滑参数如下表所示。（具体实现见代码）

\begin{table}[h]
    \centering
    \begin{tabular}{|c|c|}
    \hline
    \textbf{参数} & \textbf{值} \\ \hline
    $\alpha$ & 0.32 \\ \hline
    $\beta$ & 0.66 \\ \hline
    \end{tabular}
    \caption{双指数平滑模型的参数值}
    \end{table}

将$\alpha$和$\beta$分别代入水平（Level）更新公式和趋势（Trend）更新公式即可求出各时间点对应的水平估计和趋势估计,如下表所示。
\begin{table}[h]
    \centering
    \begin{tabular}{|c|c|c|}
    \hline
    \textbf{年份} & \textbf{水平估计 ($\ell_t$)} & \textbf{趋势估计 ($b_t$)} \\ \hline
    2015 & 603.798790 & 23.852190 \\ \hline
    2016 & 627.526523 & 23.727200 \\ \hline
    2017 & 651.129633 & 23.602865 \\ \hline
    2018 & 674.609204 & 23.479181 \\ \hline
    2019 & 697.964694 & 23.356145 \\ \hline
    2020 & 721.198425 & 23.233755 \\ \hline
    2021 & 744.309291 & 23.112005 \\ \hline
    2022 & 767.299571 & 22.990893 \\ \hline
    2023 & 790.171453 & 22.870417 \\ \hline
    \end{tabular}
    \caption{水平估计（$\ell_t$）与趋势估计（$b_t$）}
    \end{table}
    
从表中我们可以获取2023年的水平估计和趋势估计，取步长为3，4，5即可预测未来三年的教育经费投入，如下表所示。
\begin{table}[h]
    \centering
    \begin{tabular}{|c|c|}
    \hline
    \textbf{年份} & \textbf{预测教育经费 (亿元)} \\ \hline
    2026 & 812.92 \\ \hline
    2027 & 835.55 \\ \hline
    2028 & 858.07 \\ \hline
    \end{tabular}
    \caption{未来三年教育经费预测}
    \end{table}
    
\subsection{模型效果检验}
1.拟合值与残差
首先我们通过模型得到拟合值，然后计算残差。通过检查残差的性质评估模型的拟合效果。
\begin{table}[ht]
    \centering
    \begin{tabular}{|c|c|}
    \hline
    \textbf{年份} & \textbf{残差 (Residuals)} \\ \hline
    2015 & -6.28 \\ \hline
    2016 & 16.45 \\ \hline
    2017 & 7.54 \\ \hline
    2018 & 12.04 \\ \hline
    2019 & -20.23 \\ \hline
    2020 & -0.72 \\ \hline
    2021 & -35.11 \\ \hline
    2022 & -18.89 \\ \hline
    2023 & 45.16 \\ \hline
    \end{tabular}
    \caption{残差表}
    \end{table}

2.计算均方根误差、平均绝对误差、平均绝对百分比误差
\[
\text{RMSE} = \sqrt{\frac{1}{n} \sum_{t=1}^{n} (y_t - \hat{y}_t)^2}
\]

\[
\text{MAE} = \frac{1}{n} \sum_{t=1}^{n} |y_t - \hat{y}_t|
\]

\[
\text{MAPE} = \frac{1}{n} \sum_{t=1}^{n} \left|\frac{y_t - \hat{y}_t}{y_t}\right| \times 100
\]
最终结果（如下表所示）均很小，表明拟合和预测结果非常优秀。
\begin{table}[h]
    \centering
    \begin{tabular}{|c|c|}
    \hline
    \textbf{指标} & \textbf{值} \\ \hline
    RMSE & 22.49 亿元 \\ \hline
    MAE  & 18.05 亿元 \\ \hline
    MAPE & 2.50\% \\ \hline
    \end{tabular}
    \caption{模型拟合指标}
\end{table}

3.绘制Q-Q图（结果如下图所示，这里补一张qq图），数据点大致落在参考直线附近，正态性成立。

4.Ljung–Box检验。通过计算不同滞后期的自相关系数，并检验其是否显著，来判断残差序列是否存在自相关。

\[Q = n(n+2) \sum_{k=1}^{m} \frac{r_k^2}{n-k}\]

其中，$r_k$是自相关系数，n是样本大小，m是滞后期数，Q统计量用于检验是否有自相关。

将我们的数据代入公式，经检验残差不显著自相关，故可以认为我们的模型已经充分拟合，且残差类似白噪声。（具体实现见代码）

\subsection{策略分析}
为提高各阶段教育水平，结合第一问未来三年在校生人数预测，现给出策略分析如下：

1.优化教育资源配置。针对吉林省部分地区义务教育阶段教育资源相对匮乏的情况，应加大财政投入，提升教育基础设施建设，吸引人才援助落后地区，从而确保城乡教育均衡发展。
同时，应着重提升教师队伍素质，加大纪委对教务人员的监察力度，严厉打击非法补课行为，从根本上确保教师师德配位、爱护学生、有班级荣誉感，并且及时清理不作为、乱作为、嫌贫爱富、唯利是图、欺压学生的领导及教师，为吸纳的新青年教师提供良好的生态环境和晋升环境，从而促进吉林省教育事业可持续发展，长期稳定发展，留住人才。

2.考虑到吉林省高校在校生人数逐年增加的趋势，应加大财政投入，促进高校与企业的合作，推动科研成果转化落实。同时注意到中等职业教育阶段人数的稳定性，应发展与地方经济密切相关的职业教育，多向职业教育学校拨款，调整大众对职业教育的认知，加大对职业教育学校学生的监管，提高他们的平均素养。
多提供就业岗位，为区域发展提供人才支持，留住人才。

3.由于吉林省义务教育阶段在校生人数逐年减少，应加大财政投入，给青少年提供舒适的教育环境。无论是硬件设施还是软件设施，都应尽力而为、做到最好。同时多安排家访、关注学生家庭关系、亲子关系，促进家教合一，为学生提供全面的安稳的成长环境。

\end{document}

